La elección de cada técnica depende de quién posee la información y del tipo de conocimiento que se busca levantar. El conocimiento consciente es aquel que las personas recuerdan y pueden explicar al responder preguntas en entrevistas o cuestionarios, mientras que el inconsciente abarca hábitos tan cotidianos que solo se detectan al observarlos. Por su parte, el conocimiento subconsciente incluye necesidades que los usuarios aún no imaginan y que se exploran más adelante mediante prototipos o técnicas creativas. Por este motivo, el equipo combina técnicas de consulta directa con observación en campo \cite{pohl2025,ireb2026}.

\subsection{Técnicas por parte interesada}

La Tabla~\ref{tab:tecnicas} asigna a cada parte interesada la técnica de elicitación y el instrumento correspondiente.

\begin{table}[H]
  \centering\small
  \caption{Técnica de elicitación asignada a cada parte interesada.}\label{tab:tecnicas}
  \renewcommand{\arraystretch}{1.3}
  \begin{tabularx}{\textwidth}{@{}>{\raggedright\arraybackslash}p{3.0cm}>{\raggedright\arraybackslash}p{2.7cm}L>{\raggedright\arraybackslash}p{2.1cm}@{}}
    \toprule
    Parte interesada & Técnica & Justificación (fuente y tipo de conocimiento) & Instrumento \\
    \midrule
    PI-02 Encargado del gimnasio & Entrevista semiestructurada & El encargado conoce toda la operación diaria y esta técnica permite repreguntar sobre los detalles que no queden claros; recoge conocimiento consciente y detallado. & Anexo~\ref{anx:entrevista} \\
    PI-03, PI-04 y PI-05 Usuarios: estudiantes, docentes y administrativos & Cuestionario en línea & Los usuarios son numerosos, por lo que un cuestionario llega a más personas en poco tiempo y permite medir cuántos enfrentan cada problema. & Anexo~\ref{anx:cuestionario} \\
    PI-02 y PI-03 Encargado y usuarios durante un turno & Observación no participante & Muestra lo que nadie cuenta por costumbre, como quedarse después de su hora o turnarse las máquinas; es conocimiento inconsciente y ya se comprobó en la visita del 18/11. & Anexo~\ref{anx:observacion} \\
    PI-01 Responsable del gimnasio en UBU & Reunión de validación & Tiene alto poder de decisión pero poco tiempo disponible, así que una reunión corta sirve para validar el alcance y los objetivos del proyecto. & Acta de reunión \\
    PI-06 a PI-08 (según Tabla~\ref{tab:identificacion-pi}) & Entrevista breve & Se involucran en una fase posterior, una vez definido y validado el alcance principal del sistema. & Guion breve \\
    \bottomrule
  \end{tabularx}
\end{table}

\subsection{Cronograma de sesiones}

La Tabla~\ref{tab:sesiones} presenta las sesiones de elicitación. La sesión S-01 ya se realizó; las demás tienen fechas tentativas que se confirmarán con la organización.

\begin{table}[H]
  \centering\small
  \caption{Cronograma de sesiones de elicitación.}\label{tab:sesiones}
  \renewcommand{\arraystretch}{1.3}
  \begin{tabularx}{\textwidth}{@{}l>{\raggedright\arraybackslash}p{2.4cm}LL>{\raggedright\arraybackslash}p{2.6cm}@{}}
    \toprule
    Sesión & Fecha & Técnica y participantes & Artefacto esperado & Responsable \\
    \midrule
    S-01 & 18/11/2026, 08h00 a 10h00 (realizada) & Observación no participante de un turno; usuarios y encargado & Registro de observación (\url{evidencias/registro_observacion_01.md}) & E.~Monserrate, C.~Lavayen \\
    S-02 & 25/11/2026 & Entrevista semiestructurada al encargado del gimnasio & Notas de campo y transcripción anonimizada & E.~Monserrate \\
    S-03 & 26/11 al 02/12/2026 & Aplicación del cuestionario a usuarios & Respuestas anonimizadas y tabulación & C.~Lavayen, W.~Briones \\
    S-04 & 03/12/2026 & Reunión de validación con la responsable del gimnasio en UBU & Acta de reunión y alcance validado & A.~Aldas \\
    \bottomrule
  \end{tabularx}
\end{table}

\subsection{Instrumentos}

Los instrumentos se presentan en el Anexo~\ref{anx:entrevista} (guion de entrevista), el Anexo~\ref{anx:cuestionario} (cuestionario) y el Anexo~\ref{anx:observacion} (guía de observación). Todos se aplican con el consentimiento informado del Anexo~\ref{anx:consentimiento}.

La entrevista (Anexo~\ref{anx:entrevista}) aporta al ON-01, al ON-02 y al ON-03 porque detalla los pasos y demoras de la inscripción actual, las reglas de turnos y cupos por bloque horario y cómo el encargado maneja hoy las lesiones y rutinas. El cuestionario (Anexo~\ref{anx:cuestionario}) respalda el ON-01 al medir cuánto tardan los usuarios en completar la inscripción, aporta al ON-02 al identificar los horarios y la frecuencia con que asisten, y apoya el ON-03 al identificar cómo consiguen sus rutinas actualmente. Por último, la observación (Anexo~\ref{anx:observacion}) complementa el ON-02 y el ON-03 al comprobar en campo la ocupación real por turno, la falta de control al finalizar la hora y la necesidad de mayor guía durante el entrenamiento.